% Data Archiving System TRS; WBS 123456.02.01.
\section{3. User's Need Presentation}\label{users-needs}

\subsection{3.1. General User Need}\label{general-user-need}

The Data Archiving System (\abbr{DAS}) is developed as a subsystem of the EWAS Ground Segment to provide reliable storage, organisation, preservation and retrieval of mission data generated and received during EWAS operations.

The main user need addressed by the DAS is the availability of reliable and traceable historical mission data. The system shall preserve both original data received from the Ground Station Network (\abbr{GSN}) and processed data products generated by the Data Processing System (\abbr{DPS}). These data shall remain accessible to authorised users and other Ground Segment subsystems for operational reference, historical analysis, system performance evaluation and future improvement activities.

The DAS is not an independent system intended to satisfy the final user's needs directly. It supports the overall EWAS mission by providing data archiving capabilities to the Ground Segment and ensuring that relevant mission information is preserved throughout the defined retention periods.

\subsection{3.2. Intended Users}\label{intended-users}

The DAS is intended to support the following categories of users:

\begin{itemize}
    \item \textbf{Ground Segment operators:} responsible for monitoring and operating the archiving functions, supervising data ingestion and retrieval, and checking the availability and integrity of archived data.
    \item \textbf{Data analysts:} responsible for analysing historical mission data, reviewing previously detected events and evaluating the performance of the EWAS system.
    \item \textbf{Authorised crisis-management authorities:} users who may access selected archived information for historical reference, review of previous warning events and support of preparedness and response planning, subject to applicable access permissions.
    \item \textbf{Technical and maintenance personnel:} responsible for maintaining the DAS infrastructure, monitoring its technical condition, managing storage resources and performing backup and recovery activities.
\end{itemize}

In addition, selected EWAS Ground Segment subsystems may exchange data with the DAS through defined interfaces. Access to archived information shall be limited according to the operational role and authorisation level of each user or subsystem.

\subsection{3.3. Relationship with Other EWAS Products}\label{relationships-with-ewas-products}

The Data Archiving System is part of the EWAS Ground Segment and operates in cooperation with other Ground Segment subsystems. Its main role is to preserve data generated or received by these subsystems and provide access to the archived information when required.

The principal relationships between the DAS and other EWAS subsystems are described below:

\begin{itemize}
    \item \textbf{Ground Station Network (GSN):} provides raw mission data received from the spacecraft and ground station infrastructure. These data are transferred to the DAS and stored in the Raw Data Archive (\abbr{RDA}) without modification of their original content.
    \item \textbf{Data Processing System (DPS):} generates processed mission data and analytical products based on the received raw data. Selected outputs are transferred to the DAS for storage in the Processed Data Archive (\abbr{PDA}). The DPS may also retrieve archived data for historical analysis, reprocessing and evaluation activities.
    \item \textbf{Early Warning System (EWS):} generates and disseminates warning messages to authorised external users and systems based on processed information received from the DPS. Generated warning messages and associated dissemination information, including timestamps, dissemination status and event identifiers, are transferred to the DAS for archival purposes and subsequent historical review.
\end{itemize}

The DAS does not perform the primary processing of mission data, generate early warning decisions or control spacecraft and ground station operations. These responsibilities remain within the relevant EWAS subsystems.

\subsection{3.4. Identified Interfaces}\label{identified-interfaces}

The DAS requires interfaces with the relevant Ground Segment elements to receive, store and provide access to mission data. The principal expected interfaces are:

\begin{itemize}
    \item \textbf{GSN--DAS interface:} transfer of raw mission data and associated acquisition metadata from the Ground Station Network to the Raw Data Archive.
    \item \textbf{DPS--DAS interface:} transfer of processed data products, event records, analytical results and associated processing metadata from the Data Processing System to the Processed Data Archive. The interface may also support retrieval of archived data by the DPS.
    \item \textbf{DAS--EWS interface:} enables the transfer and archival of generated warning messages and associated dissemination information, including timestamps, dissemination status and relevant event identifiers. This interface ensures the preservation of warning history and supports subsequent review, traceability and evaluation of the EWAS warning process.
\end{itemize}

The detailed technical characteristics, data exchange formats, communication protocols and access mechanisms of these interfaces shall be defined during the subsequent design and interface specification activities.

\subsection{3.5. Expected Benefits}\label{expected-benefits}

The implementation of the DAS provides the EWAS project with a persistent and organised source of historical mission information. The availability of archived raw and processed data supports the investigation of previous events, verification of processing results, identification of potential errors and evaluation of system performance.

Furthermore, the preservation of relevant metadata, processing history and event classifications enables historical datasets to be used in future analysis and in the training, validation and evaluation of detection and assessment algorithms. The DAS provides the necessary data preservation and retrieval capabilities for these activities but does not itself perform algorithm training or operational data processing.
